\section{Zadanie 1 (Mateusz Wojaczek)}
Niech $X$ będzie zmienną losową o rozkładzie normalnym $N(\mu,\sigma^2)$.
Wyznaczyć wszystkie momenty centralne zmiennej $X$, tj.
\[
\mathbb{E}[(X-\mu)^k] \quad \text{dla } k \in \mathbb{N}.
\]

\subsection*{Rozwiązanie}
Niech
\[
Y \sim N(0,1).
\]
Wtedy
\[
X = \mu + \sigma Y.
\]

Zatem
\[
\mathbb{E}[(X-\mu)^k] = \mathbb{E}[(\sigma Y)^k] = \sigma^k \mathbb{E}[Y^k].
\]

Obliczamy więc momenty zmiennej $Y$.

\subsubsection*{Obliczenie $\mathbb{E}(Y^k)$}

Z definicji:
\[
\mathbb{E}[Y^k]
= \frac{1}{\sqrt{2\pi}}
\int_{0}^{\infty} y^k e^{-\frac{y^2}{2}} \, dy.
\]

Oznaczmy:
\[
I_k = \int_{0}^{\infty} y^k e^{-\frac{y^2}{2}} \, dy.
\]

Obliczamy całkę $I_k$ przez całkowanie przez części:
\[
I_k
= \int_{0}^{\infty} y^{k-1} \cdot y e^{-\frac{y^2}{2}} \, dy
\]

\[
= \left[ y^{k-1} \left( -e^{-\frac{y^2}{2}} \right) \right]_{0}^{\infty}
+ \int_{0}^{\infty} (k-1) y^{k-2} e^{-\frac{y^2}{2}} \, dy.
\]

Wyraz brzegowy znika, zatem otrzymujemy rekurencję:
\[
I_k = (k-1) I_{k-2}.
\]

\subsubsection*{Wartości początkowe}

\[
I_0 = \int_{0}^{\infty} e^{-\frac{y^2}{2}} \, dy = \sqrt{2\pi},
\]

\[
I_1 = 0,
\]
ponieważ jest to wartość oczekiwana zmiennej $Y \sim N(0,1)$.

\subsubsection*{Postać ogólna}

Dla $k$ parzystego:
\[
I_{2k} = \prod_{i=1}^{k} (2i-1)\,\sqrt{2\pi}.
\]

Dla $k$ nieparzystego:
\[
I_{2k+1} = 0.
\]

Stąd:
\[
\mathbb{E}[Y^k] =
\begin{cases}
0, & \text{gdy $k$ jest nieparzyste}, \\
\displaystyle \prod_{i=1}^{k-1} (k-2i), & \text{gdy $k$ jest parzyste}.
\end{cases}
\]

\subsubsection*{Wniosek końcowy}

\[
\mathbb{E}[(X-\mu)^k] = \sigma^k \mathbb{E}[Y^k],
\]
czyli:
\[
\mathbb{E}[(X-\mu)^k] =
\begin{cases}
0, & \text{gdy $k$ jest nieparzyste}, \\
\sigma^k \displaystyle \prod_{i=1}^{\frac{k}{2}-1} (k-2i), & \text{gdy $k$ jest parzyste}.
\end{cases}
\]

\subsection*{Odpowiedź}
\[
\mathbb{E}[(X-\mu)^k] =
\begin{cases}
0, & \text{gdy $k$ jest nieparzyste}, \\
\sigma^k \displaystyle \prod_{i=1}^{\frac{k}{2}-1} (k-2i), & \text{gdy $k$ jest parzyste}.
\end{cases}
\]
